% !TEX root = ../main.tex
\chapter{PCSEL 的耦合波理论（CWT）：最小模型}
\label{ch:cwt-min}

\section*{学习目标}
\begin{itemize}
  \item 理解为什么在能带图之后必须引入耦合波理论（Coupled-Wave Theory, CWT），以及它在 PCSEL 研究链条中的准确位置。
  \item 从 Maxwell 方程与场展开出发，建立最小 CWT 模型，明确每一步近似究竟丢掉了什么物理。
  \item 学会把四波 CWT 基底改写成 $\Gamma$ 点 little-group 的 irrep 基底，从而把“耦合矩阵”直接连到“模标注、偏振与 BIC 判断”。
  \item 学会把耦合系数、辐射损耗、材料增益和模竞争放进同一组低维方程中阅读，而不是把它们拆成互不相干的名词。
\end{itemize}

\section*{先修知识}
建议已掌握 \cref{ch:maxwell,ch:gamma,ch:feedback,ch:threshold} 中关于开放 Maxwell 问题、Gamma 点候选模、面内耦合与阈值条件的内容。

\section*{本章路线图}
\ChapterModelLevel{L1→L2}
PCSEL 的面发射来自周期散射将面内导模耦合到垂直辐射通道；阈值问题要求把增益和损耗写入同一动力学框架。CWT 提供介于能带图和全波仿真之间的低维模型。本章先说明能带图的不足，再从场展开和投影出发建立最小模型，解释矩阵中各类项的物理来源，并把四波基底改写成 $\Gamma$ 点 irrep 基底，使模式标签、偏振选择和辐射选择规则进入同一矩阵语言。最后说明其在设计流程中的用法与边界。

\noindent\textbf{记号约定（避免符号歧义）}：本章与导波相关的 $\beta$ 默认都指 $\betaslab$（slab 面内传播常数）；自发辐射耦合因子统一写作 $\betasp$，仅在动态与噪声章节使用。

\begin{RigorousBox}
除非另作声明，本章的时间因子约定与全书一致，均采用 $\ee^{-\ii\omega t}$。因此后文所有等效损耗、增益和复频率符号，均应与 \cref{ch:maxwell,ch:math-appendix} 的频域约定联读，而不能另行切换正负号解释。
\end{RigorousBox}

\section{为什么在能带图之后还需要 CWT}
对初学者来说，能带图往往已经很有吸引力。它告诉我们哪些 Gamma 点模值得关注，也告诉我们模之间可能存在简并、分裂与光锥耦合。但是，能带图主要回答的是“有哪些候选模”，而不是“这些候选模为什么会互相竞争、又为什么某种几何扰动会特别有效”。后一类问题需要把注意力从“单个本征模的频率”转移到“若干主导平面波分量之间如何相互耦合”。

CWT 的作用正是在此处引入。它不替代 Maxwell 方程，也不否定全波仿真的必要性；它是在已知目标模主要面内波分量后，将这些分量投影为低维耦合系统。这样，几何参数的作用可通过耦合系数、损耗项和增益项逐项解释。

\begin{IntuitionBox}
PWE 和能带图主要给出候选模列表；RCWA、FDTD、FEM 用于高保真验证。CWT 位于二者之间，用于解释候选模如何因少数散射通道重组，以及这些通道增强或削弱后阈值、偏振和远场的变化趋势。
\end{IntuitionBox}

\section{最小 CWT 要解决的物理问题是什么}
在方格晶格 PCSEL 的最典型工作区间里，Gamma 点附近的目标模常常主要由四个等价的面内导模分量组成：沿 $+x$、$-x$、$+y$、$-y$ 方向传播的基本波。上一章已经解释过，周期折射率调制通过倒格矢把这些分量连接起来，并进一步把其中的某些组合折叠进光锥，形成垂直辐射。因而，若只关心 Gamma 点附近最核心的反馈网络，一个自然的最小基底就是这四个波分量。

于是我们想做的事就变得很明确：

第一，把三维矢量 Maxwell 场写成“纵向模剖面”乘以“少数几个面内平面波”的组合；

第二，把周期结构导致的相互散射压缩为若干耦合系数；

第三，把辐射损耗、内部损耗和增益作为等效非厄米项写进同一套方程；

第四，通过求这个低维系统的本征值和本征向量，判断哪一类线性组合更可能成为实际候选激射模。

若这四步成立，所得矩阵不是任意经验矩阵，而是由 Maxwell 问题投影得到的物理压缩模型。

\begin{figure}[htbp]
  \centering
  \begin{tikzpicture}[x=1cm,y=1cm,>=Latex]
    \begin{scope}
      \draw[fill=blue!8,rounded corners] (-1.6,-1.2) rectangle (1.6,1.2);
      \draw[->,thick] (0,0)--(1.15,0) node[right] {$R_x$};
      \draw[->,thick] (0,0)--(-1.15,0) node[left] {$S_x$};
      \draw[->,thick] (0,0)--(0,0.95) node[above] {$R_y$};
      \draw[->,thick] (0,0)--(0,-0.95) node[below] {$S_y$};
      \node at (0,-1.7){实空间/面内传播图像};
    \end{scope}
    \begin{scope}[xshift=7.2cm]
      \draw[->] (-1.8,0)--(1.8,0) node[right] {$k_x$};
      \draw[->] (0,-1.8)--(0,1.8) node[above] {$k_y$};
      \fill[blue!70] (1.2,0) circle (0.07) node[below] {$+\beta\hat{\bm x}$};
      \fill[blue!70] (-1.2,0) circle (0.07) node[below] {$-\beta\hat{\bm x}$};
      \fill[blue!70] (0,1.2) circle (0.07) node[right] {$+\beta\hat{\bm y}$};
      \fill[blue!70] (0,-1.2) circle (0.07) node[right] {$-\beta\hat{\bm y}$};
      \draw[<->,red!75,thick] (1.2,0)--(-1.2,0);
      \draw[<->,red!75,thick] (0,1.2)--(0,-1.2);
      \draw[<->,orange!80!black,thick] (1.2,0)--(0,1.2);
      \draw[<->,orange!80!black,thick] (1.2,0)--(0,-1.2);
      \draw[<->,orange!80!black,thick] (-1.2,0)--(0,1.2);
      \draw[<->,orange!80!black,thick] (-1.2,0)--(0,-1.2);
      \node[red!75] at (0,0.23) {$\kappa_1:\,\bm G_x$};
      \node[orange!80!black] at (0.86,0.78) {$\kappa_2:\,\bm G_x\!\pm\!\bm G_y$};
      \node[fill=white,inner sep=1pt] at (0,-0.55) {$A_1,\;B_1,\;E$};
      \node at (0,-2.0){倒空间中的四波与主导散射路径};
    \end{scope}
  \end{tikzpicture}
  \caption{最小四波 CWT 图像。左图强调四支基本波在面内传播；右图强调它们在倒空间中如何通过主导倒格矢散射彼此连接。对 PCSEL 而言，后续的偏振、BIC 与阈值排序都建立在这类最小闭合散射网络之上。}
  \label{fig:cwt_four_wave_map}
\end{figure}

\section{从 Maxwell 方程到场展开：四波近似从哪里来}
为了不让推导显得像突然写下一个猜测矩阵，我们从背景波导出发。设器件在没有横向周期调制时，已经支持一个主要的纵向导模剖面 $\bm{e}_0(z)$。当加入二维周期折射率扰动后，我们在 Gamma 点附近只保留四个最主要的面内波分量，写成
\begin{equation}
\E(\rho,z,t)
\approx
\ee^{-\ii\omega t}
\bm{e}_0(z)
\sum_{m=1}^{4} a_m(\rho)
\ee^{\ii\bm{\beta}_m\cdot \rho},
\label{eq:cwt_ansatz}
\end{equation}
其中 $\rho=(x,y)$，四个传播向量取为
\[
\bm{\beta}_1=(\beta,0),\quad
\bm{\beta}_2=(-\beta,0),\quad
\bm{\beta}_3=(0,\beta),\quad
\bm{\beta}_4=(0,-\beta).
\]
这里的 $a_m(\rho)$ 是缓慢变化的包络，表示每个基本波分量在大尺度上的权重；$\bm{e}_0(z)$ 则把纵向层栈的信息压缩进一个主要导模剖面。之所以允许这种写法，是因为我们关注的是 Gamma 点附近、频率接近且相互强耦合的一小族波，而不是整个倒格空间中无穷多个分量。

接下来，把介电常数写成背景项与周期扰动之和：
\begin{equation}
\varepsilon(\rho,z)=\varepsilon_0(z)+\Delta\varepsilon(\rho,z),
\qquad
\Delta\varepsilon(\rho+\bm{R},z)=\Delta\varepsilon(\rho,z).
\label{eq:cwt_eps_split}
\end{equation}
再将 $\Delta\varepsilon$ 作二维傅里叶展开，
\begin{equation}
\Delta\varepsilon(\rho,z)
=
\sum_{\gvec}
\Delta\varepsilon_{\gvec}(z)\ee^{\ii\gvec\cdot\rho},
\label{eq:cwt_delta_eps_fourier}
\end{equation}
便得到不同 $\bm{\beta}_m$ 之间的散射通道。把 \cref{eq:cwt_ansatz,eq:cwt_delta_eps_fourier} 相乘时，每一个基本波 $a_n\ee^{\ii\bm\beta_n\cdot\rho}$ 都会生成一组
\[
\Delta\varepsilon_{\gvec}(z)\,a_n(\rho)\,
\ee^{\ii(\bm\beta_n+\gvec)\cdot\rho}
\]
项。于是，投影到第 $m$ 个基底后，真正留下来的耦合必须满足
\begin{equation}
\bm\beta_m \approx \bm\beta_n+\gvec
\qquad\Longleftrightarrow\qquad
\gvec \approx \bm\beta_m-\bm\beta_n.
\label{eq:cwt_selection_rule}
\end{equation}
这就是 CWT 里“耦合矩阵元来自哪一个 Fourier 分量”的最小选择定则。只保留“近共振”的倒格矢耦合后，四个包络之间便不再独立，而是由少数几个主导傅里叶分量彼此连接。

\section{投影步骤到底在做什么}
形式上，将 \cref{eq:cwt_ansatz} 代回频域 Maxwell 方程，再用背景导模剖面 $\bm{e}_0(z)$ 和四个平面波基函数作为测试函数做加权投影。这个步骤等价于限定系统只在四个主导波分量张成的子空间中演化，并求出 Maxwell 算符在该子空间内的矩阵表示。

进行投影时，通常还要用到两个额外近似。

第一是慢变包络近似（slowly varying envelope approximation）。它要求 $a_m(\rho)$ 相对载波相位因子变化得慢，从而二阶横向导数和远离共振的高频混频项可以忽略。这一步使方程从“完整波动问题”降到“包络动力学问题”。
若想把“变化得慢”说得更定量一些，可把包络变化长度记为 $L_{\mathrm{env}}$，则最小要求是
\[
L_{\mathrm{env}}\gg a,
\qquad
\big|\nabla_\rho^2 a_m\big|\ll \beta^2 |a_m|.
\]
它的物理含义是：一个晶格常数尺度上，主要由快速 Bloch 相位承载振荡；包络只在更大的器件尺度上记录边缘、增益和热斑导致的慢变化。对典型 PCSEL 而言，若晶格常数 $a\approx\SI{280}{\nano\meter}$、有限阵列直径 $D\approx\SIrange{100}{300}{\micro\meter}$，则 $D/a\approx 360$--$1100$，这正是最小 CWT 在器件中央区通常仍有教学价值的原因。也正因如此，一旦进入边缘急剧截断、局部热斑或强非均匀注入区，最小 CWT 的慢变包络假设就会先变得可疑。

第二是主导耦合截断。周期结构的傅里叶分量原则上可以把一个波分量耦合到许多其他分量，但真正主导 Gamma 点物理的往往是少数几个满足近布拉格条件的通道。最小 CWT 正是主动舍弃那些次要通道，以换取可解释性。

把这个“截断”说得更公式化一点：若本征频率附近存在许多可能的平面波分量，则只有同时满足下列两类条件的通道才值得进入最小模型：
\begin{enumerate}
  \item \textbf{动量近匹配}：满足 \cref{eq:cwt_selection_rule}；
  \item \textbf{失谐足够小}：对应未耦合导模频率差 $|\omega_m^{(0)}-\omega_n^{(0)}|$ 不大于主导耦合系数数量级。
\end{enumerate}
对方格晶格、Gamma 点附近、二阶 Bragg 工作的最小教材模型来说，$\pm x$ 与 $\pm y$ 这四支波恰好同时满足这两类要求，因此构成自然的四波闭合子空间；更远的倒格矢通道则因频率失谐或纵向重叠太弱而被先行舍弃。这一截断的定量充分条件可粗略写成：主导傅里叶分量幅度满足 $|\Delta\varepsilon_{\mathbf{G}}|/\bar{\varepsilon}\ll 1$（折射率对比度足够小），同时次近邻布拉格通道的失谐远大于主导耦合系数 $|\kappa|$。若这两个条件不满足，则应升级到更多波分量的展开或直接转入 RCWA/FDTD 全波方法。

完成这些步骤后，就得到线性耦合方程
\begin{equation}
\frac{\dd}{\dd t}\bm{a}
=
\left(\ii\bm{\Omega}-\bm{\Gamma}+\bm{G}\right)\bm{a},
\qquad
\bm{a}=
\begin{bmatrix}
R_x & S_x & R_y & S_y
\end{bmatrix}^{\mathsf T},
\label{eq:cwt_matrix}
\end{equation}
其中 $R_x,S_x,R_y,S_y$ 分别表示四个基本波分量的包络振幅。\cref{eq:cwt_matrix} 看起来像一个普通矩阵方程，但每个矩阵块背后都有清楚的来源和职责，下面要逐项解释。

若进一步假设包络按单一复频率偏移演化，
\begin{equation}
\bm a(t)=\bm u\,\ee^{-\ii\delta\omega t},
\label{eq:cwt_time_harmonic_ansatz}
\end{equation}
则在本书统一采用的时间因子 $\ee^{-\ii\omega t}$ 约定下，时间导数满足
\[
\frac{\dd}{\dd t}\rightarrow -\ii\delta\omega.
\]
把 \cref{eq:cwt_time_harmonic_ansatz} 代回 \cref{eq:cwt_matrix}，便得到
\begin{equation}
\left[
-\ii\delta\omega\,\bm I
-\left(\ii\bm{\Omega}-\bm{\Gamma}+\bm{G}\right)
\right]\bm u=0.
\label{eq:cwt_frequency_domain_bridge}
\end{equation}
\cref{eq:cwt_frequency_domain_bridge} 给出了从时间演化方程到频域本征值问题的最小转换。后文高阶或三维 CWT 中写成 $(\delta+\ii\alpha)\bm V=\bm C\bm V$，只是对这一步进行不同记号重排与归一化，并非另一套独立理论。

若把投影步骤写得再直白一点，它其实就是以下三步。
\begin{enumerate}
  \item 先选定低维基底 $\{\bm e_0(z)\ee^{\ii\bm\beta_m\cdot\rho}\}$；
  \item 再把完整 Maxwell 算符夹在这组基底左右，形成矩阵元
  \[
  M_{mn}\sim
  \left\langle
  \bm e_0\ee^{\ii\bm\beta_m\cdot\rho},
  \hat{\mathcal L}
  \bm e_0\ee^{\ii\bm\beta_n\cdot\rho}
  \right\rangle;
  \]
  \item 最后只保留满足 \cref{eq:cwt_selection_rule} 的主导矩阵元，把其余远离共振的通道当作高阶修正。
\end{enumerate}
因此，CWT 不是猜测一个 $4\times4$ 矩阵，而是把完整 Maxwell 问题投影到预先选定且物理上主导的有限维共振子空间。

\begin{RigorousBox}
严格地说，\cref{eq:cwt_matrix} 不是从完整 Maxwell 方程“精确等价变形”而来，而是从完整问题向低维共振子空间做投影之后的近似闭合系统。它的可信度取决于两件事：一是被丢弃的基底是否真的次要，二是工作点是否确实靠近所选共振子空间。脱离这两个条件，CWT 只能作为启发图像，而不能作为定量工具。
\end{RigorousBox}

\cref{fig:ch09_cwt_signal_flow}给出了CWT 信号流的完整可视化：四个行波基底通过倒格矢散射形成闭合耦合网络，其中$\kappa_1$描述反向散射（1D DFB 型），$\kappa_2$描述90° 转角散射（2D PCSEL特有），辐射损耗通过光锥耦合项体现，材料增益作为对角项作用于所有基底。

\begin{figure}[htbp]
  \centering
  % !TEX root = ../main.tex
% Figure content only: inserted inside an outer figure environment in ch09
\begin{tikzpicture}[x=1cm,y=1cm]
  \tikzset{
    wavenode/.style={circle,draw=blue!70,thick,fill=blue!10,minimum size=1.05cm,font=\small},
    coup1/.style={<->,>=Stealth,thick,red!75},
    coup2/.style={<->,>=Stealth,thick,orange!85!black},
    loss/.style={->,>=Stealth,thick,dashed,red!60},
    gain/.style={->,>=Stealth,thick,green!50!black}
  }

  % A
  \begin{scope}[xshift=0cm,yshift=3.5cm]
    \node[font=\bfseries] at (0,2.2) {(a) 四波基底};
    \node[wavenode] (rx) at (-1.8,0.9) {$R_x$};
    \node[wavenode] (sx) at (1.8,0.9) {$S_x$};
    \node[wavenode] (ry) at (-1.8,-0.9) {$R_y$};
    \node[wavenode] (sy) at (1.8,-0.9) {$S_y$};
    \draw[coup1] (rx) -- node[above,font=\scriptsize] {$\kappa_1$} (sx);
    \draw[coup1] (ry) -- node[below,font=\scriptsize] {$\kappa_1$} (sy);
    \draw[coup2] (rx) -- node[left,font=\scriptsize] {$\kappa_2$} (ry);
    \draw[coup2] (sx) -- node[right,font=\scriptsize] {$\kappa_2$} (sy);
    \draw[coup2] (rx) -- (sy);
    \draw[coup2] (sx) -- (ry);
    \draw[loss] (rx) -- ++(0,0.9) node[above,font=\scriptsize] {$\gamma_{\rm rad}$};
    \draw[loss] (sx) -- ++(0,0.9);
    \draw[loss] (ry) -- ++(0,-0.9) node[below,font=\scriptsize] {$\gamma_{\rm rad}$};
    \draw[loss] (sy) -- ++(0,-0.9);
    \draw[gain] (-3.3,0.9) -- node[above,font=\scriptsize] {$g_{\rm mod}$} (rx);
    \draw[gain] (3.3,0.9) -- (sx);
    \draw[gain] (-3.3,-0.9) -- node[below,font=\scriptsize] {$g_{\rm mod}$} (ry);
    \draw[gain] (3.3,-0.9) -- (sy);
  \end{scope}

  % B
  \begin{scope}[xshift=-3.8cm,yshift=-1.2cm]
    \node[font=\bfseries] at (2.2,1.7) {(b) 矩阵形式};
    \node[align=center,font=\small] at (2.2,-0.1) {
      $\displaystyle
      \frac{\dd}{\dd t}
      \begin{bmatrix}
      R_x\\ S_x\\ R_y\\ S_y
      \end{bmatrix}
      =
      \left(\ii\bm{\Omega}-\bm{\Gamma}+\bm{G}\right)
      \begin{bmatrix}
      R_x\\ S_x\\ R_y\\ S_y
      \end{bmatrix}
      $
    };
    \node[draw=black,rounded corners,fill=gray!6,align=left,font=\footnotesize,text width=4.8cm] at (2.2,-2.1) {
      \textbf{读法：}
      $\bm{\Omega}$ 负责保守耦合，
      $\bm{\Gamma}$ 负责辐射/吸收损耗，
      $\bm{G}$ 负责模增益。
    };
  \end{scope}

  % C
  \begin{scope}[xshift=4.8cm,yshift=-1.2cm]
    \node[font=\bfseries] at (1.8,1.7) {(c) 信号流解释};
    \node[draw=blue!70,rounded corners,fill=blue!6,align=left,font=\footnotesize,text width=4.6cm] at (1.8,-0.5) {
      \textbf{$\kappa_1$：}反向散射，类似 1D DFB 中的前后向反馈。\\[3pt]
      \textbf{$\kappa_2$：}转角散射，把 $x$ 向和 $y$ 向波连接起来，是 2D PCSEL 的关键。\\[3pt]
      \textbf{$\gamma_{\rm rad}$：}通向光锥的辐射泄漏。\\[3pt]
      \textbf{$g_{\rm mod}$：}有源区对四波基底的等效增益注入。
    };
  \end{scope}
\end{tikzpicture}

  \caption{CWT信号流与耦合网络。(a)四波ansatz的实空间图像：正向行波$R_x$、反向行波$S_x$、侧向行波$R_y$和$S_y$在面内传播；(b)倒空间中的耦合路径：红色双向箭头表示由倒格矢$\gvec_x$驱动的反向散射（耦合强度$\kappa_1$），橙色双向箭头表示由$\gvec_x\pm\gvec_y$驱动的90°转角散射（耦合强度$\kappa_2$）；(c)完整的信号流图：每个节点代表一个行波基底的复振幅，边上的标记表示对应的耦合机制（DFB反馈、转角散射、辐射泄漏、材料增益），闭环增益条件要求绕任意闭合回路一周的总增益等于总损耗。}
  \label{fig:ch09_cwt_signal_flow}
\end{figure}

\cref{fig:ch09_cwt_signal_flow} 将耦合矩阵元素对应到散射网络。$\kappa_1$ 对应传统 1D DFB 激光器的反向散射机制，$\kappa_2$ 对应 2D PCSEL 中 $x$ 与 $y$ 方向波分量之间的转角散射；后者使面内反馈具有二维特征，并参与表面发射通道的形成。

\section{三类矩阵项各自代表什么}
先看 $\bm{\Omega}$。它负责相位演化和相互散射，也就是通常所说的“保守耦合”部分。若只讨论被动、无损、完全封闭的近似系统，$\bm{\Omega}$ 的本征值主要决定模的频率位置和简并如何被打破。几何参数如孔半径、填充比和晶格常数，首先往往就是通过改变 $\bm{\Omega}$ 中的耦合系数来改写模的排列。

再看 $\bm{\Gamma}$。它并不只是一项“把数值减掉”的阻尼常数，而是所有损耗通道在该低维基底中的压缩表示。对 PCSEL 来说，其中至少要区分两类来源：一类是垂直辐射损耗，它来自模与自由空间通道的耦合；另一类是内部损耗，包括自由载流子吸收、掺杂层吸收、金属邻近损耗和粗糙散射等。也就是说，$\bm{\Gamma}$ 本身就已经把“开放系统”信息带进了模型。

最后是 $\bm{G}$。如果只讨论冷腔，这一项可以先不写；但一旦要谈阈值，$\bm{G}$ 就不可回避。它描述的不是抽象的”有源”二字，而是有源区增益在这组波分量基底中的投影。最简单的教学写法是
\[
\bm{G}=g_{\mathrm{mod}}\bm{I},
\]
即先假设四个主波共享同一个均匀模增益。若增益分布近似均匀，这确实是一个自然起点；
但若注入明显非均匀，或不同波分量对应的场在横向上有不同重叠，那么 $\bm{G}$ 的对角元就会分裂，
进一步表现出模式选择性。

\begin{RigorousBox}
$g_{\mathrm{mod}}$ 的具体形式取决于量子阱材料、注入条件和工作温度。
在 \cref{ch:qw-gain} 中，材料增益被展开为
$g_{\mathrm{mat}}(N,T,\hbar\omega)$；
经由约束因子 $\Gamma_{\mathrm{opt}}$ 投影后得到 $g_{\mathrm{mod}}$，
再代入 $\bm G=g_{\mathrm{mod}}\bm I$（或更一般的对角分裂形式）。
若跳过这一步而把 $g_{\mathrm{mod}}$ 当作自由拟合参数，则 CWT 的可解释性会大幅下降，因为它把“有源区物理”变成了纯粹的数字调参。
\end{RigorousBox}

于是，\cref{eq:cwt_matrix} 的本质就很清楚了：
\begin{itemize}
  \item $\bm{\Omega}$ 决定“哪些组合容易形成共振”；
  \item $\bm{\Gamma}$ 决定“这些组合会以多快的速度泄露或耗散”；
  \item $\bm{G}$ 决定“有源区是否优先把增益送给其中某些组合”。
\end{itemize}
真正的激射候选模，是这三者共同作用下的本征态，而不是只由其中某一项决定。

\subsection{先说清归一化：重叠积分和耦合系数的数值可比性}
在解析推导和数值仿真中，读者会频繁看到“场归一化”“模归一化”“包络振幅”这些说法。它们背后对应的并不是唯一标准。

在被动本征问题中，本征场常常只确定到一个乘法常数，因此场幅度本身并不直接对应实际功率；可比较的是频率、损耗率、归一化后的重叠积分等量。在有源问题中，若把模展开代入耦合模方程，包络振幅的定义又取决于采用的是能量归一化、功率归一化还是某种方便解析的尺度。

因此，本书对归一化采取一个务实原则：只要某个公式中出现重叠积分、模增益或耦合系数，就必须在附近说明归一化方式属于哪一类；如果没有说明，就不能把不同章节或不同软件的数值直接拿来比较。

本书后续会频繁使用两种最常见的归一化方式，先在此给出可计算的标准写法：
\begin{align}
\text{能量归一化（被动本征模常用）：}
&&\int_V \varepsilon_r(\rvec)\,|\E_n(\rvec)|^2\,\mathrm{d}V &= 1, \\
\text{功率归一化（导波/端口问题常用）：}
&&\frac{1}{2}\int_S \operatorname{Re}(\E_n\times\Hf_n^\ast)\cdot\hat{z}\,\mathrm{d}A &= 1.
\end{align}
其中 $S$ 为与传播方向正交的截面；对开放 slab 结构若需统计上下两个辐射通道，应明确 $S$ 是单侧截面还是上下两侧合并截面，以免功率计数出现两倍差异。两种归一化得到的模幅 $a_m$ 单位不同，耦合系数 $\kappa$ 的量纲也随之不同。跨章节或跨软件比较 $\kappa$ 数值时，必须先确认双方使用同一种归一化。


\section{耦合系数为什么是几何和场重叠的压缩表达}
最关键的非对角元来自周期扰动 $\Delta\varepsilon$ 对不同平面波分量的散射。把投影过程写开，典型的耦合系数具有如下结构：
\begin{equation}
\kappa_{mn}
\propto
\int_V
\Delta\varepsilon(\rho,z)
\big[\bm{e}_0^{\ast}(z)\cdot\bm{e}_0(z)\big]
\ee^{-\ii(\bm{\beta}_m-\bm{\beta}_n)\cdot\rho}
\dd V.
\label{eq:cwt_kappa}
\end{equation}
\cref{eq:cwt_kappa} 表明耦合系数可由扰动 Fourier 分量和纵向场重叠估计，而不只是拟合参数。它给出三点结论。

第一，耦合并不是由几何尺寸单独决定的，而是由“几何的傅里叶分量”和“纵向模场重叠”共同决定。因此，改变外延层厚度或波导层折射率，也会通过 $\bm{e}_0(z)$ 改写面内耦合强度。

第二，不同的几何扰动并不会平均地影响所有耦合通道。某种孔形改变可能主要放大 $x$ 向和 $y$ 向之间的耦合，另一种扰动则可能主要增强进入光锥的散射。也正因为如此，PCSEL 设计不能只说“增加扰动”或“减小扰动”，而必须问清楚“增强的是哪一类傅里叶分量”。

第三，如果某个耦合系数同时参与面内反馈和垂直辐射，那么它天然就对应设计中的权衡。因为增强它，可能一边在提升反馈，另一边也在提升辐射损耗。

把 \cref{eq:cwt_kappa} 和前面的选择定则结合起来，还能直接读出 $\kappa_1,\kappa_2$ 为什么恰好由少数 Fourier 分量给定。对最小四波模型而言，沿 $\pm x$ 的前后向耦合主要由
\[
\bm\beta_1-\bm\beta_2\approx \left(\frac{2\pi}{a},0\right)
\]
对应的一族 Fourier 分量控制；沿 $x/y$ 两组波之间的交叉耦合则主要由
\[
\bm\beta_1-\bm\beta_3,\quad
\bm\beta_1-\bm\beta_4
\]
所对应的 Fourier 分量控制。于是，$\kappa_1$、$\kappa_2$ 并不是两个凭经验命名的常数，而是“沿不同倒格矢方向的几何 Fourier 分量，经过同一纵向模场剖面重加权后得到的压缩参数”。

为了避免“从 Maxwell 方程突然跳到一个耦合矩阵”显得过快，这里至少把一个矩阵元显式写出来。设方格晶格最主要的倒格矢为
\[
\bm G_x=\left(\frac{2\pi}{a},0\right),
\]
并关注沿 $\pm x$ 传播的两支基本波之间的耦合，即 $\bm\beta_1=(\beta,0)$ 与 $\bm\beta_2=(-\beta,0)$。把周期扰动展开为
\[
\Delta\varepsilon(\rho,z)=\sum_{\gvec}\Delta\varepsilon_{\gvec}(z)\ee^{\ii\gvec\cdot\rho},
\]
则由 \cref{eq:cwt_kappa} 得
\begin{align}
\kappa_{12}
&\propto
\int_V
\sum_{\gvec}\Delta\varepsilon_{\gvec}(z)\ee^{\ii\gvec\cdot\rho}
\big[\bm e_0^\ast(z)\cdot \bm e_0(z)\big]
\ee^{-\ii(\bm\beta_1-\bm\beta_2)\cdot\rho}\,\dd V
\notag\\
&=
\sum_{\gvec}
\int \Delta\varepsilon_{\gvec}(z)\big|\bm e_0(z)\big|^2\,\dd z
\int_{A_c}
\ee^{\ii[\gvec-(2\beta,0)]\cdot\rho}\,\dd^2\rho.
\label{eq:cwt_matrix_element_example}
\end{align}
在近布拉格条件 $2\beta\approx 2\pi/a$ 下，面积积分只保留 $\gvec=\bm G_x$ 这一项，于是
\begin{equation}
\kappa_{12}\equiv \kappa_1
\propto
A_c
\int \Delta\varepsilon_{\bm G_x}(z)\big|\bm e_0(z)\big|^2\,\dd z.
\label{eq:cwt_kappa1_example}
\end{equation}
\cref{eq:cwt_kappa1_example} 说明矩阵元来自\textbf{特定倒格矢 Fourier 分量}与\textbf{纵向模场重叠}的乘积，而不是抽象常数。其他矩阵元类似，只是对应不同的 $\bm\beta_m-\bm\beta_n$ 与不同的主导倒格矢。

\begin{ExampleBox}
对方格晶格圆形空气孔（半径 $r$、晶格常数 $a$、背景折射率 $n_b\approx3.4$），主导倒格矢 $\bm G_x=(2\pi/a,0)$ 对应的截面 Fourier 分量可解析估算为
\[
\Delta\varepsilon_{\bm G_x}\approx 2\left(\frac{r}{a}\right) J_1\!\left(\frac{2\pi r}{a}\right) \cdot(n_{\rm air}^2-n_b^2)
\approx -2\left(\frac{r}{a}\right)J_1\!\left(\frac{2\pi r}{a}\right)\cdot n_b^2,
\]
其中 $J_1$ 是一阶 Bessel 函数。对典型占空比 $r/a\approx0.3$，$J_1(2\pi\times0.3)\approx0.40$，因此
\[
|\Delta\varepsilon_{\bm G_x}|\approx 2\times0.3\times0.40\times11.6\approx2.8.
\]
再乘以纵向重叠积分（\textbf{[本书算例]}中光子晶体层常贡献约 10--20\% 纵向模能量），有效耦合系数 $\kappa_1$ 的典型量级约为 $\kappa_1/(\omega/c) \sim 10^{-2}$--$10^{-1}$，对应波矢修正 $\delta k \sim |\kappa_1|/v_g\sim1$--$10\,\mathrm{cm}^{-1}$。这应读作“数量级一致性”而非通用拟合律；与已报道 CWT/实验量级对照可见于\cite{liang2012square,liang2013triangular,yoshida2023}。
\end{ExampleBox}

\section{四个基本波如何重新组合成实际候选模}
四波基底真正有价值的地方，不在于记住 $R_x,S_x,R_y,S_y$ 四个名字，而在于理解它们会重新组合成更适合讨论对称性的超模。以方格晶格为例，常见的线性组合包括“全同相”的亮模、“沿某方向相消”的暗模，以及在 $x$、$y$ 两个方向上具有不同奇偶性的组合。不同线性组合与辐射通道的重叠不同，因而辐射损耗也不同。

这一步正好把前面几章连起来：\cref{ch:gamma} 告诉我们 Gamma 点为什么容易产生简并和模式分裂，\cref{ch:feedback} 解释了哪些散射通道打开了垂直辐射，而在 CWT 里，这些叙述第一次统一成了一个低维本征问题。求解 \cref{eq:cwt_matrix} 的本征向量，我们得到的就是“由四个基本波相干叠加形成的候选超模”；求其本征值，则能读出频率位置和净增长或净衰减趋势。

如果把 \cref{eq:cwt_matrix} 写成本征值形式
\begin{equation}
\left(\ii\bm{\Omega}-\bm{\Gamma}+\bm{G}\right)\bm{u}_j
=
\lambda_j\bm{u}_j,
\label{eq:cwt_eig}
\end{equation}
那么 $\bm{u}_j$ 就给出第 $j$ 个候选超模在四个基本波基底上的权重，而 $\lambda_j$ 的实部则表征该超模在小信号近似下是净增长还是净衰减。

\section{把四波基底改写成 \texorpdfstring{$C_{4v}$}{C4v} 的 irrep 基底}
为了使 CWT 成为可复用的教材工具，不能只说明四个基本波可形成若干组合，还应把这些组合写成 \cref{ch:polarization} little-group 框架中的不可约表示。对方格晶格，$\Gamma$ 点 little group 为 $C_{4v}$。在最小四波子空间
\[
\bm{a}=
\begin{bmatrix}
R_x & S_x & R_y & S_y
\end{bmatrix}^{\mathsf T}
\]
中，一个特别有用的 irrep 基底可取为
\begin{align}
\bm{u}_{A_1}
&=
\frac{1}{2}
\begin{bmatrix}
1\\1\\1\\1
\end{bmatrix},
\label{eq:cwt_A1_basis}
\\
\bm{u}_{B_1}
&=
\frac{1}{2}
\begin{bmatrix}
1\\1\\-1\\-1
\end{bmatrix},
\label{eq:cwt_B1_basis}
\\
\bm{u}_{E,x}
&=
\frac{1}{\sqrt{2}}
\begin{bmatrix}
1\\-1\\0\\0
\end{bmatrix},
\qquad
\bm{u}_{E,y}
=
\frac{1}{\sqrt{2}}
\begin{bmatrix}
0\\0\\1\\-1
\end{bmatrix}.
\label{eq:cwt_E_basis}
\end{align}
其中 $\bm{u}_{A_1}$ 是全同相 singlet，$\bm{u}_{B_1}$ 描述 $x$ 向与 $y$ 向波对之间的相对反相，而 $\bm{u}_{E,x},\bm{u}_{E,y}$ 则张成二维 $E$ 双重态，对应最直接的 $x$-like 与 $y$-like 偏振支路。

把这四个列向量拼成变换矩阵
\begin{equation}
U = \frac{1}{\sqrt{2}}
\begin{bmatrix}
\frac{1}{\sqrt{2}} & \frac{1}{\sqrt{2}} & 1 & 0 \\
\frac{1}{\sqrt{2}} & \frac{1}{\sqrt{2}} & -1 & 0 \\
\frac{1}{\sqrt{2}} & -\frac{1}{\sqrt{2}} & 0 & 1 \\
\frac{1}{\sqrt{2}} & -\frac{1}{\sqrt{2}} & 0 & -1
\end{bmatrix},
\label{eq:cwt_U_matrix}
\end{equation}
则对理想方格晶格（$\kappa_1,\kappa_2$ 均为实数），CWT 矩阵 $\tilde H$ 在 irrep 基底 $\bm{b}=U^\dagger\bm{a}$ 下化为块对角形式：$A_1$ 与 $B_1$ 各对应一个 $1\times1$ 块，$E$ 双重态对应一个 $2\times2$ 块（在 $E$ 子空间中仍可能存在内部耦合，但与 $A_1,B_1$ 块解耦）。

这一写法说明，四波 CWT 中应关注的不是单个 $R_x$ 或 $S_y$ 分量，而是这些分量在 little-group 意义下如何组织成 $A_1$、$B_1$ 和 $E$ family。于是：
\begin{itemize}
  \item 模式标注不再依赖“看图感觉像哪种场型”，而可以直接由 irrep 给出；
  \item $E$ 双重态天然与 \cref{ch:polarization} 中的双偏振/法向辐射通道对应；
  \item $A_1$ 和 $B_1$ 则天然提示“这是 symmetry-distinct singlet”，其辐射与混合规则必须和 $E$ 分开看。
\end{itemize}

\begin{RigorousBox}
\cref{eq:cwt_A1_basis,eq:cwt_B1_basis,eq:cwt_E_basis} 是\textbf{最小四波子空间中的 irrep 分类}，而不是整个 Maxwell 问题所有 irrep 的完整列表。在这个最小子空间里，$A_2$ 和 $B_2$ 没有显式出现；若加入更多 Fourier 分量、更多纵向模或更完整矢量分量，它们可以在更高维投影中出现。
\end{RigorousBox}

这件事最好再说得更物理一些。最小四波基底只保留
\[
(\pm\beta,0)\hat{\bm y},\qquad (0,\pm\beta)\hat{\bm x}
\]
这四支 TE-like 基本波，因此它张成的是一个四维表示。对该表示做 character reduction，
恰好只分解出 $A_1\oplus B_1\oplus E$。要得到 $A_2/B_2$，至少还需要再引入额外的场自由度，
例如更高阶 Fourier 分量、不同纵向场型，或使镜面对称下的奇偶结构足以支撑新的 singlet。
因此 $A_2/B_2$ 的缺席不是“被我们忘了写”，而是最小四波子空间本身就不承载它们。

\section{这套 irrep 基底怎样直接连接偏振、辐射与 BIC 判断}
现在可以把 \cref{ch:polarization} 和上一章的分类工具直接落到 CWT 矩阵里。对方格晶格，法向自由空间辐射通道在最小 in-plane 分类下属于 $E$。因此，在保持 $C_{4v}$ 对称的前提下：
\begin{itemize}
  \item $E$ 双重态最直接地与法向辐射通道匹配，因此往往承担“可法向辐射偏振支路”的角色；
  \item $A_1$ 与 $B_1$ 这类 singlet 即使存在，也不会与 $E$ 通道发生对称性允许的直接混合；
  \item 若某一 singlet family 在 exact $\Gamma$ 点呈现 direct normal-radiation suppression，就应优先把它放进 symmetry-protected dark family / BIC 候选的语境中理解。
\end{itemize}

于是，CWT 的矩阵语言第一次真正和“irrep 标注”“偏振选择”“BIC 初筛”闭合起来：求 \cref{eq:cwt_eig} 的本征向量之后，读者不仅能说“哪个模本征值最大”，还能说“这个本征模在 irrep 意义下属于哪一类，它在对称性上是否允许法向辐射，它与哪一支偏振通道最直接相关”。

\section{对称性破缺怎样在 CWT 矩阵里体现为 irrep 分裂}
当结构保持理想方格对称时，CWT 矩阵在 irrep 基底中往往接近分块形式：$A_1$、$B_1$ 和 $E$ family 彼此分离。引入小对称性破缺后，例如圆孔变椭圆、双晶格位移或外延/电极造成的有效各向异性，这种分块结构会被重新组织。

对常见偏振工程问题，关键步骤是：二维 $E$ 双重态在降群后分裂成两支非简并 singlet，使 $x$-like 与 $y$-like 偏振不再等价。同时，某些原本 direct radiation 受限的 family 可能因降群允许混合而获得有限辐射损耗。因此，在 CWT 中，“偏振钉扎”“模式分裂”“quasi-BIC 打开”可视为\textbf{同一耦合矩阵在不同对称性下的不同投影结果}。

\begin{ExampleBox}
若方格晶格圆孔的 $C_{4v}$ 对称被轻微椭圆孔降到只保留两条镜面对称的子群，则原先的 $E$ 双重态会沿椭圆主轴分裂成两支非简并偏振模。在 CWT 矩阵里，这表现为 $\bm{u}_{E,x}$ 与 $\bm{u}_{E,y}$ 所在子块不再共享同一本征值。若其中一支同时更强地接近法向辐射通道，另一支则更接近 dark family，就会同时观察到偏振选择、损耗分裂与远场主瓣定向。
\end{ExampleBox}

\section{为什么二维八波 CWT 仍然不够：它不是错误，而是压缩得太狠}
已有的二维八波 CWT 对 PCSEL 入门非常重要，因为它第一次把“四个基本波 $+$ 若干高阶中继波”这件事说清楚了。但教材里应进一步说明：二维模型之所以会在真实器件中失效，不是因为它方向错了，而是因为它压缩了三个在真实 PCSEL 中会重新变重要的自由度。
\begin{enumerate}
  \item \textbf{纵向场自由度被压掉了。} 二维模型默认不同波矢分量共享相同的垂直场型，而真实器件中的基本波、辐射波和高阶波在 $z$ 方向上往往看见的是不同的波导环境；
  \item \textbf{高阶 Fourier 分量被截断得过狠。} 对圆孔、弱调制、近对称结构，这种截断常可接受；但对深刻蚀、非对称孔形、多孔单胞和三晶格结构，高阶 Fourier 分量对二维耦合本身就可能是主导项；
  \item \textbf{垂直辐射被过度经验化。} 若只用一个拟合辐射常数去代表所有法向耦合，就很难可靠处理上下出光不对称、远场整形与结构微扰引起的辐射重排。
\end{enumerate}
因此，“从二维 CWT 升级到三维 CWT”的含义不是单纯增加模型复杂度，而是\textbf{显式纳入真实 PCSEL 中被二维模型压缩、但会重新影响结果的通道}。

\section{三维 CWT 能算什么：一页速览}
上述三点压缩，在 \cref{ch:cwt-3d} 中会被逐项恢复：Liang--Noda 三维 CWT 从矢量波动方程出发，让辐射波与高阶波各自携带真实的纵向场型并用 Green 函数解析求解，最终仍归结为一个 $4\times4$ 本征问题，其中每个矩阵元都由孔形、填充比、晶格常数和外延层栈解析确定，原则上不需要经验拟合参数。

在进入完整推导之前，先记住几组数字（均来自 \cite{liang2012square}），它们解释了为什么这套模型在设计流程中不可替代：
\begin{itemize}
  \item \textbf{速度}：一组参数下求解 $4\times4$ 本征问题约需 $1$ 秒，而同等精度的 3D FDTD 约需 $4$ 小时；整条能带（约 100 个 $\kvec$ 点）三维 CWT 约 $2$ 分钟，3D FDTD 需要数百小时。
  \item \textbf{精度}：模式频率预测与 3D FDTD 吻合到 $<1\%$。
  \item \textbf{物理输出}：对方格晶格圆孔，模式 A/B 在精确 $\Gamma$ 点辐射常数为零（即 \cref{ch:bic} 的 symmetry-protected BIC）；换成直角等腰三角形孔后 $\alpha_{\mathrm{rad}}$ 约为 $50$--$200\,\mathrm{cm}^{-1}$，$Q$ 从无穷大变为有限可设计。
  \item \textbf{收敛代价}：辐射常数收敛需要把平面波截断取到 $D\geq 10$（即 $441$ 个分量），而频率收敛 $D\geq 5$ 已经足够——两类输出的收敛门槛并不相同。
\end{itemize}

完整推导、有限尺寸包络方程和非 $\Gamma$ 点能带扩展见 \cref{ch:cwt-3d}；从 CWT 本征矢量提取远场偏振（方位角、椭圆率与 Stokes 参量）的实用公式已收入 \cref{ch:math-appendix}。本章余下部分回到最小模型能够独立回答的问题：阈值条件、模式竞争与设计语言。


\section{CWT 怎样表达阈值条件}
现在把 \cref{ch:threshold} 的阈值条件接进来。如果仅讨论光学阈值而不求阈值电流，那么在 CWT 里最直接的说法是：当某个本征值的实部恰好穿过零时，该超模达到线性阈值。也就是说，阈值条件可写成
\begin{equation}
\max_j \operatorname{Re}(\lambda_j)=0.
\label{eq:cwt_threshold}
\end{equation}
这个表达比“$g_{\mathrm{th}}=\alpha_{\mathrm{rad}}+\alpha_{\mathrm{int}}$”更适合描述 PCSEL 的模竞争，因为它允许不同超模具有不同的辐射损耗、内部损耗投影和增益重叠。CWT 并不推翻 \cref{ch:threshold} 的阈值公式，而是把单模阈值条件推广为多模耦合后的超模本征值问题。

若把第 $j$ 支超模的本征值显式拆成
\begin{equation}
\lambda_j
=
v_{\mathrm{en},j}
\bigl[
\Gamma_{\mathrm{opt},j}g_{\mathrm{mat}}
-\alpha_{\mathrm{rad},j}
-\alpha_{\mathrm{int},j}
\bigr]
+ \ii\,\Delta\omega_j,
\label{eq:cwt_threshold_bridge}
\end{equation}
那么 \cref{eq:cwt_threshold} 就立刻退化回
\begin{equation}
\Gamma_{\mathrm{opt},j}g_{\mathrm{mat}}
=
\alpha_{\mathrm{rad},j}
+\alpha_{\mathrm{int},j}
\qquad
(\text{对最先起振的 }j).
\label{eq:cwt_threshold_bridge_scalar}
\end{equation}
这正是 \cref{ch:threshold} 中单模阈值平衡式在 CWT 超模语言下的写法。
区别只在于：Ch07 把这些量写成单模参数，CWT 则允许它们随超模 family 改变。

但也必须保持严格：\cref{eq:cwt_threshold} 给出的仍然只是光学线性阈值。若要把它变成阈值电流，仍然必须知道 $\bm{G}$ 如何随注入电流、载流子分布和温度变化。这一部分必须由 \cref{ch:qw-gain,ch:electrical,ch:eto-coupling} 提供。

\section{从线性阈值到模式竞争动力学}
阈值条件 $\max_j\operatorname{Re}\lambda_j=0$ 给出的只是"哪一支超模最先达到光学阈值"。但真实器件一旦跨过阈值，就不再是单一阈值问题，而是一个多模竞争的动力学问题。这是因为：各超模 family 通常拥有不同的净增益和不同的辐射/重叠特性，一旦某支模率先达到阈值，其光场就会通过增益饱和反过来压制其他模的净增益，而饱和本身又与注入、空间烧孔和非均匀电流分布交织在一起。

CWT 的本征值方程在这一层级已经不能直接使用，因为 \cref{eq:cwt_matrix} 中的 $\bm{a}$ 默认为固定背景场中的线性包络，而阈值以上的激光场要求增益、载流子密度和光场三者满足自洽条件。更准确的动力学语言是耦合到多模速率方程：
\begin{equation}
\frac{\dd S_j}{\dd t}
=
\left[\Gamma_{\mathrm{opt},j}g_{\mathrm{mat}}(N,T)
-\alpha_{\mathrm{rad},j}
-\alpha_{\mathrm{int},j}\right]S_j
+\beta_{\mathrm{sp}}R_{\mathrm{sp},j},
\qquad j=1,\ldots,4,
\label{eq:cwt_multimode_rate}
\end{equation}
其中 $S_j$ 是第 $j$ 支超模的光子数，$\Gamma_{\mathrm{opt},j}g_{\mathrm{mat}}$ 中的 $g_{\mathrm{mat}}$ 不再是常数，而是通过载流子速率方程与所有光子数共同决定：
\begin{equation}
\frac{\dd N}{\dd t}
=
\frac{I}{eV_{\mathrm{act}}}
-\left[R_{\mathrm{rad}}+R_{\mathrm{nr}}\right]
-\sum_j \frac{\Gamma_{\mathrm{opt},j}}{v_{\mathrm{en}}}\,
g_{\mathrm{mat}}(N,T)\,S_j.
\label{eq:cwt_carrier_rate}
\end{equation}

在这条方程链中，每个超模的 $\Gamma_{\mathrm{opt},j}$、$\alpha_{\mathrm{rad},j}$ 和 $\alpha_{\mathrm{int},j}$ 来自 CWT 本征值分析；注入电流 $I$、载流子密度 $N$ 和温度 $T$ 由电注入、热反馈和增益模型给出。因此，CWT 提供模式竞争的线性输入，阈值以上哪一支模占优还需由速率方程或自洽动力学模型决定。

从教学编排上说，本章之所以把动力学问题留到这里才提，是因为读者必须先知道"各超模的净损耗和重叠因子从哪来"（CWT 的职责），才能追问"一旦跨过阈值，这些参数如何被速率方程重新分配"（\cref{ch:dynamics} 的职责）。若跳过这一步，直接把 $\max_j\operatorname{Re}\lambda_j=0$ 当成最终器件结论，就会在阈值和动态稳定性之间留下一个不该有的概念空白。

\begin{RigorousBox}
\cref{eq:cwt_multimode_rate} 中的 $\Gamma_{\mathrm{opt},j}$、$\alpha_{\mathrm{rad},j}$ 和 $\alpha_{\mathrm{int},j}$ 均来自 CWT 本征值分解，它们已经是给定结构几何和外延层栈下的"冷冻"参数。
真正的动态竞争发生在阈值以上：先起振的模通过增益饱和改变 $g_{\mathrm{mat}}(N,T)$，而 $g_{\mathrm{mat}}$ 的改变又反过来重新定义各模的净增益——这个反馈回路只能用时域或稳态自洽方程描述，而不能用 CWT 的线性本征值直接读出。
\end{RigorousBox}

\section{一个最典型的设计问题：几何扰动到底在帮你什么}
很多设计讨论会说“加一个小扰动以获得单模”或者“调孔形以改善偏振”，但这类说法只有落到 CWT 的语言里才真正可检验。设某个小扰动主要改变了 $x$ 向与 $y$ 向基本波之间的耦合，而几乎不改变与光锥的直接耦合，那么它更可能首先体现在模式分裂和偏振选择上，而未必显著降低阈值。相反，如果某个扰动主要增强了与辐射通道相关的矩阵元，那么它可能更直接改变辐射损耗和提取效率。

\begin{ExampleBox}
考虑两种孔形修改方案。方案 A 主要改变晶格的二阶傅里叶分量，使 $x$ 向与 $y$ 向波分量的交叉耦合显著增强；方案 B 主要增强与零面内动量辐射通道相关的傅里叶分量。用 CWT 语言看，方案 A 更容易引起模式分裂和偏振选择，方案 B 更容易改写辐射损耗与输出耦合。若目标是降低阈值并维持单峰远场，两者的优先级可能完全不同。
\end{ExampleBox}

这个例子说明，CWT 的价值不只在于计算速度，还在于把设计动作分解为可逐项检查的物理效果。

\subsection{三晶格如何直接进入 CWT：Fourier 放大因子与 \texorpdfstring{$180^\circ$}{180 deg} 耦合增强}
若要把多孔单胞或三晶格结构真正纳入教材，而不是只把它当成“经验上更强反馈”的例子，最好的方式是直接在 Fourier 系数层面写清楚。设单晶格光子晶体的介电扰动写成
\begin{equation}
  \Delta\varepsilon_{1}(x,y)
  =
  \sum_{m,n}F_{m,n}\exp\!\left[\ii\frac{2\pi}{a}(mx+ny)\right].
\end{equation}
若在同一晶胞中再叠加两个相对位移分别为 $(d_{1x},d_{1y})$ 和 $(d_{2x},d_{2y})$ 的嵌套晶格，则三晶格总扰动满足
\begin{equation}
  \Delta\varepsilon_{\mathrm{3L}}(x,y)
  =
  \sum_{m,n}A_{m,n}F_{m,n}\exp\!\left[\ii\frac{2\pi}{a}(mx+ny)\right],
\end{equation}
其中放大因子为
\begin{equation}
  A_{m,n}
  =
  1+\exp\!\left[-\ii\frac{2\pi}{a}(md_{1x}+nd_{1y})\right]
  +\exp\!\left[-\ii\frac{2\pi}{a}(md_{2x}+nd_{2y})\right].
\label{eq:trilattice_amplification}
\end{equation}

对正方晶格 PCSEL 而言，四个基本波之间最关键的是 $180^\circ$ 反向耦合，因此首先应关注 $(m,n)=(2,0)$ 与 $(0,2)$：
\begin{align}
  A_{2,0} &= 1+\exp\!\left(-\ii\frac{4\pi}{a}d_{1x}\right)+\exp\!\left(-\ii\frac{4\pi}{a}d_{2x}\right), \\
  A_{0,2} &= 1+\exp\!\left(-\ii\frac{4\pi}{a}d_{1y}\right)+\exp\!\left(-\ii\frac{4\pi}{a}d_{2y}\right),
\end{align}
因此有
\begin{equation}
  0\le |A_{2,0}|,\ |A_{0,2}|\le 3.
\label{eq:trilattice_amplification_bound}
\end{equation}
这条结果表明，三晶格的优势不在于单纯增加孔数，而在于可将关键二阶 Fourier 分量从相消调到三倍相干增强，从而定向增强 $180^\circ$ 反馈通道 \cite{liang2013triangular,wang2024}。

对设计而言，这还有两层重要含义。第一，若把 $A_{2,0}$ 或 $A_{0,2}$ 调到接近零，就会系统性削弱反向反馈；第二，若把二者同时调到接近 3，则可显著增强面内光反馈、降低水平泄露，并为更小尺寸谐振腔创造条件。但边界仍需保留：增强 $180^\circ$ 耦合首先改善的是冷腔反馈与模式损耗排序，并不自动等于器件阈值电流最低。真实阈值仍要继续经过 $\bm C_{\mathrm{rad}}$、内部吸收、量子阱重叠、注入与热回写的再筛选。

\section{CWT 到底能算什么，不能算什么}
如果用得恰当，最小 CWT 至少能做五件非常有价值的事。

第一，它能给出候选超模的构成，告诉我们目标模主要由哪些基本波叠加而成。

第二，它能预测某些几何参数变化会沿哪条物理通道起作用，是更多地改频率、改损耗，还是改模式分裂。

第三，它能用很低的计算代价做灵敏度扫描，为 RCWA、FDTD、FEM 缩小参数搜索空间。

第四，它能在冷腔或简化增益模型下，快速比较几个候选模的相对阈值趋势。

第五，它能把“偏振工程”“远场工程”“模式抑制”放进同一种耦合语言下理解，而不是把这些主题拆成彼此独立的经验技巧。

但它同样有明确的不能之处。

它不能严格替代全波方法去计算开放三维结构的精确辐射损耗；不能在没有外延、电学和热学模型的前提下直接给出阈值电流；不能可靠处理远离所选共振子空间的大尺度模重组；也不能自动处理有限尺寸边缘、注入不均匀和强热漂移引起的复杂模式变形。

\begin{PitfallBox}
把 CWT 输出的“目标模本征值最大”直接当成“器件一定单模工作”的证据，是非常常见的误判。CWT 最多告诉你在所选基底、所选近似和所选增益模型下，某个超模线性起振更有优势。真实器件还要接受有限尺寸、热漂移、注入不均匀和非线性增益压缩的再筛选。
\end{PitfallBox}

\section{CWT 在整本书工作流中的准确位置}
把本书的建模链条回顾一遍，就能看清 CWT 的职责边界。

在 \cref{ch:bloch,ch:gamma,ch:feedback} 中，我们建立的是模式分类和面发射的物理图像；在 \cref{ch:threshold,ch:polarization,ch:bic} 中，我们把阈值、偏振和远场放进开放系统视角；从本章开始，则进入“半解析模型”层。这里的目标不是追求一次性得出最终器件结论，而是建立一套足够清楚的中间语言，使后续 RCWA、FDTD、FEM 以及器件级电热模型之间能互相对话。

因此，最合理的研究流程通常不是“CWT 或全波二选一”，而是：
\begin{enumerate}
  \item 用 PWE 或经验直觉找到 Gamma 点附近的候选模族；
  \item 用 CWT 理解主导耦合通道和参数灵敏度；
  \item 用 RCWA 精修开放周期单胞的辐射与谱线行为；
  \item 用 FEM/FDTD 验证有限尺寸器件模式；
  \item 再把这些结果送入外延、电注入和热耦合模型，判断器件级阈值与稳定工作区。
\end{enumerate}

因此，CWT 在 PCSEL 研究中的作用不在于最高精度，而在于充当跨模型的中间描述。

\begin{quote}
能带章节回答的是“哪些 $\Gamma$ 点或近 $\Gamma$ 点带边值得进入候选名单”；CWT 章节回答的是“这些候选模为何会因为 $180^\circ$ 反馈、二维高阶中继耦合与法向辐射通道而被重新排序”。二维 PWEM、三维 CWT 与后续 RCWA/FDTD/FEM 不是彼此替代关系，而是同一个 PCSEL 问题在不同层级上的投影。
\end{quote}

\begin{RigorousBox}
\cref{ch:cwt-3d} 中的三维 CWT 推导（\cref{sec:3d-cwt-derivation}）基于 Liang 等人
\cite{liang2012square,liang2013triangular} 针对方格和三角晶格的工作。
该框架已被应用于三晶格 PCSEL 设计 \cite{wang2024}，其中\textbf{[文献值]}耦合系数
$\kappa_{\mathrm{1D}}$ 的计算直接指导了三个子晶格间距的优化，将
$180^{\circ}$ 反馈耦合强度相比单晶格提升约三倍，从而显著降低阈值电流。
这说明三维 CWT 可作为孔形、子晶格位移等工程参数优化的半解析工具。
\end{RigorousBox}

\begin{CommonPitfallBox}{CWT 无限周期结果外推到有限阵列}
\textbf{错误理解：}“CWT 给出某候选超模最优，有限阵列器件一定保持同排序。”  
\textbf{正确做法：}把 CWT 作为候选筛选与灵敏度解释层，再用有限阵列全波验证尺寸效应。  
\textbf{最小验收：}至少报告 3 个阵列尺寸下的主模/竞争模排序、远场主瓣角与旁瓣变化，确认不存在尺寸引发的排序翻转。
\end{CommonPitfallBox}

\section{2024--2026 前沿补充：CWT 的极化与晶格泛化正在变成可复用工具链}
传统教学常把 CWT 先绑定到“方格晶格 + 近 $\Gamma$ + 主偏振”的最小场景。近年的进展表明，这个边界正在被系统扩展：例如三角晶格、{TM} 偏振下的二维耦合波理论已经给出可复核的解析系数与模频预测，并与有限器件计算形成了同向证据\cite{robinson2025tmcwt}。这对本章最直接的意义是：CWT 不再只是“解释器”，而是在更广晶格族上可用于前端参数筛选的半解析求解器。

与此同时，快速设计工具的系统比较工作也指出：CWT 及其迭代扩展在 2D 设计空间中仍然具备很高性价比，但其输出必须经由更高保真方法做边界校准\cite{lang2025rapidpcseltools}。因此，本章建议把 CWT 的工程使用口径固定为：
\begin{itemize}
  \item 用于快速识别“哪类耦合通道最可能改写模式排序”；
  \item 用于给出参数扫描的方向性而非最终器件结论；
  \item 在进入发布级结论前，必须与 RCWA/FDTD/FEM 的有限尺寸结果闭环互证。
\end{itemize}

\section*{本章小结}
最小 CWT 的核心思想是：在 Gamma 点附近选取少数主导平面波分量，
把完整 Maxwell 问题投影到一个低维耦合子空间中，再把保守耦合、开放损耗和增益
同时写进同一组矩阵方程。这样做的收益是把几何、外延和有源效应压缩成可解释的耦合参数；
代价则是模型只在所选子空间和近似条件内可信。本章进一步说明，这个低维子空间不应只按
“哪四个波分量”来读，还应按 little-group 的 irrep 来读。对方格晶格最小四波模型，
$A_1$、$B_1$ 与 $E$ family 的分解，正是把模式标注、偏振选择和 BIC 初筛接到同一个矩阵语言上的关键一步。
对 PCSEL 来说，CWT 既不是能带图的重复，也不是全波求解器的廉价替身，
而是连接物理图像、对称性分类与定量设计的中间层。

\section*{练习题}
\begin{enumerate}
  \item \basicq 从 \cref{eq:cwt_ansatz} 出发，说明为什么 Gamma 点附近最自然的最小基底是四个基本波，而不是任意选取四个平面波。

  \item \basicq 验证 \cref{eq:cwt_A1_basis,eq:cwt_B1_basis,eq:cwt_E_basis} 给出的四个向量在最小四波子空间中构成一组正交基，并说明为什么它们适合作为 $C_{4v}$ 的最小 irrep 基底。

  \item \midq 结合 \cref{eq:cwt_kappa}，讨论改变孔半径、改变刻蚀深度和改变纵向层栈厚度三种动作分别会怎样影响耦合系数。

  \item \midq 说明为什么在保持 $C_{4v}$ 的前提下，$E$ family 与 $A_1/B_1$ family 不能直接混合；再说明一旦降群，这种结论为什么可能被改写。

  \item \midq 解释为什么 \cref{eq:cwt_threshold} 给出的是光学线性阈值，而不是器件级阈值电流。

  \item \hardq 设计一个你认为“最小 CWT 会失效”的场景，并指出失效是因为缺少哪些基底或物理通道。
\end{enumerate}

\section*{建议继续阅读}
\begin{itemize}
  \item 下一章讨论三维修正后，可回到本章核对二维有效模型在何种 PCSEL 结构上会系统性失效。

  \item 如需与标准光子晶体教材对照，可参考 \cite{joannopoulos2008,sakoda2005} 中关于平面波耦合、对称性和模式投影的相关章节。

  \item PCSEL 专用 CWT 的代表性建立过程可参考 \cite{liang2012square,liang2013triangular}。
\end{itemize}
